% !TEX program = xelatex
\documentclass[a4paper,11pt,openany,fontset=none]{ctexrep}

\newcommand{\reporttitle}{全球头部另类资产管理人的方法论：\newline 中国个人投资者可迁移的公开市场框架}
\newcommand{\reportsubtitle}{复制投资委员会，不复制收购、杠杆与专属交易}
\newcommand{\reportkicker}{STRATEGIC INVESTMENT RESEARCH}
\newcommand{\reportscope}{GLOBAL ALTERNATIVES \textbar\ CHINA PUBLIC MARKETS}
\newcommand{\reportdate}{2026年8月13日}
\newcommand{\reportdatacutoff}{资料截至 2026年8月13日}
\newcommand{\reporttype}{研究框架证据备忘录}
\newcommand{\reportauthor}{AStock Research Agent}
\newcommand{\reporthouseview}{有可借用的研究方法，但它不是“私募买什么就跟什么”。个人可借用的是：窄能力圈、现金流与负债承保、可验证的经营改善、情景估值、分阶段复核和预先写下的失效条件。控制权、杠杆收购（LBO）、私募信贷发起、保险资金与专属交易是机构能力，应明确排除。}
\newcommand{\reportquality}{以 Blackstone、KKR、Apollo、Brookfield、Carlyle、Bain Capital 与 TPG 的官方披露，以及中国证监会、沪深交易所现行规则为主。管理人材料可证明其描述的流程，不能独立证明其超额收益可被个人复制；未使用任何用户账户、实时行情或个券数据。}
\newcommand{\reportdisclaimer}{本报告基于公开资料整理，仅供研究和教育使用，不构成任何证券买卖、产品适当性、收益或执行建议。}

% Reusable institutional-style preamble for XeLaTeX AStock research reports.
% Usage: define \reporttitle and \reportdate, then % Reusable institutional-style preamble for XeLaTeX AStock research reports.
% Usage: define \reporttitle and \reportdate, then % Reusable institutional-style preamble for XeLaTeX AStock research reports.
% Usage: define \reporttitle and \reportdate, then \input{.agents/templates/preamble.tex}.
% Design benchmark: JPM first-page dashboard, BlackRock house-view framing,
% Vanguard scenario discipline, and AB/JPM chartbook exhibit structure.

\providecommand{\reporttitle}{AStock Research Report}
\providecommand{\reportsubtitle}{House View, Evidence, Valuation and Risk}
\providecommand{\reportkicker}{INSTITUTIONAL EQUITY RESEARCH}
\providecommand{\reportscope}{CHINA A-SHARES}
\providecommand{\reportdate}{\today}
\providecommand{\reportdatacutoff}{Data cutoff not specified}
\providecommand{\reporttype}{Research note}
\providecommand{\reportauthor}{AStock Research Agent}
\providecommand{\reporthouseview}{State the house view in one decision-useful paragraph.}
\providecommand{\reportquality}{Source quality not specified.}
\providecommand{\reportdisclaimer}{本报告基于公开资料整理，不构成任何证券买卖建议。}

% ============ 页面布局 ============
\usepackage[top=1.8cm, bottom=1.9cm, left=1.8cm, right=1.8cm]{geometry}
\setlength{\parskip}{0.3em}
\setlength{\parindent}{2em}
\renewcommand{\baselinestretch}{1.12}
\raggedbottom

% ============ Fonts ============
% Prefer the installed Noto CJK family when available (for Linux builds).
% Otherwise retain the original MacTeX typography. This avoids platform
% probing via shell escape and keeps the source portable across both engines.
\IfFontExistsTF{Noto Serif CJK SC}{
  \setCJKmainfont{Noto Serif CJK SC}
  \setCJKsansfont{Noto Sans CJK SC}
  \setCJKmonofont{Noto Sans CJK SC}
  \setmainfont{Liberation Serif}
  \newcommand{\kaishu}{\CJKfontspec{Noto Serif CJK SC}}
  \newcommand{\heiti}{\CJKfontspec{Noto Sans CJK SC}}
  \newcommand{\songti}{\CJKfontspec{Noto Serif CJK SC}}
}{
  \setCJKmainfont[BoldFont={Heiti SC}]{STSong}
  \setCJKsansfont{Heiti SC}
  \setCJKmonofont{Heiti SC}
  \setmainfont{Times New Roman}
  \newcommand{\kaishu}{\CJKfontspec{STKaiti}}
  \newcommand{\heiti}{\CJKfontspec{Heiti SC}}
  \newcommand{\songti}{\CJKfontspec{STSong}}
}

% ============ 颜色（IB Navy/Grey palette） ============
\usepackage{xcolor}
\definecolor{navy}{HTML}{003366}
\definecolor{deepnavy}{HTML}{0B1F33}
\definecolor{steelgrey}{HTML}{4A5568}
\definecolor{mutedgrey}{HTML}{718096}
\definecolor{lightgrey}{HTML}{F7FAFC}
\definecolor{panelgrey}{HTML}{EDF2F7}
\definecolor{riskred}{HTML}{C53030}
\definecolor{riskamber}{HTML}{D69E2E}
\definecolor{riskgreen}{HTML}{38A169}
\definecolor{riskblue}{HTML}{2B6CB0}
\definecolor{accentblue}{HTML}{2B6CB0}
\definecolor{nvgreen}{HTML}{76B900}
\definecolor{rulegrey}{HTML}{CBD5E0}

% ============ 表格 ============
\usepackage{booktabs}
\usepackage{longtable}
\usepackage{tabularx}
\usepackage{multirow}
\usepackage{array}
\usepackage{makecell}
\usepackage{colortbl}
\renewcommand{\arraystretch}{1.12}
\newcolumntype{Y}{>{\centering\arraybackslash}X}
\newcolumntype{L}[1]{>{\raggedright\arraybackslash}p{#1}}
\newcolumntype{C}[1]{>{\centering\arraybackslash}p{#1}}
\newcolumntype{R}[1]{>{\raggedleft\arraybackslash}p{#1}}

% ============ 图形 ============
\usepackage{graphicx}
\usepackage{tikz}
\usetikzlibrary{shapes,arrows.meta,positioning,calc,backgrounds,fit,decorations.pathreplacing}
\usepackage{pgfplots}
\usepgfplotslibrary{polar}
\pgfplotsset{compat=1.18}
\usepackage[edges]{forest}

% ============ 页眉页脚（报告标题和日期需在main.tex中覆盖） ============
\usepackage{fancyhdr}
\pagestyle{fancy}
\fancyhf{}
\fancyhead[L]{\scriptsize\color{steelgrey} \reporttitle}
\fancyhead[R]{\scriptsize\color{steelgrey} \reportdate}
\fancyfoot[C]{\small\color{steelgrey} \thepage}
\fancyfoot[R]{\tiny\color{steelgrey} \reportdisclaimer}
\renewcommand{\headrulewidth}{0.4pt}
\renewcommand{\footrulewidth}{0.2pt}

% ============ 章节格式 ============
\usepackage{titlesec}
\titleformat{\chapter}[block]
  {\normalfont\LARGE\bfseries\color{navy}}
  {\thechapter}{0.9em}{}
\titlespacing*{\chapter}{0pt}{0.4em}{1.0em}
\titleformat{\section}
  {\normalfont\large\bfseries\color{navy}}
  {\thesection}{1em}{}
\titleformat{\subsection}
  {\normalfont\normalsize\bfseries\color{steelgrey}}
  {\thesubsection}{1em}{}

% ============ 超链接 ============
\usepackage{hyperref}
\hypersetup{
  colorlinks=true,
  linkcolor=navy,
  urlcolor=navy,
  citecolor=navy,
}

% ============ 其他工具 ============
\usepackage{enumitem}
\setlist{nosep, leftmargin=2em}
\IfFileExists{fontawesome5.sty}{\usepackage{fontawesome5}}{}
% ============ Class B 全局兜底（xeCJK + 缺字 fallback 阻断）：
% 当 fontawesome5 因 Times New Roman 缺字 / \FA 命名空间失败时，
% 禁止输出 aExclamationTriangle / aCloud 裸文本，一律降级为 amssymb 原生符号。
% 注：amssymb 在 L112 已加载，此处不重复。所有 \def 在本段落。
\def\faExclamationTriangle{\ensuremath{\textcolor{riskred}{\blacktriangle}}\kern-0.25em}
\def\faCloud{\ensuremath{\clubsuit}\kern-0.1em}
\def\faYenSign{\ensuremath{\yen}}
\def\faIndustry{\ensuremath{\S}}
\def\faCogs{\ensuremath{\ast}}
\def\faFlagCheckered{\ensuremath{\blacktriangleright}}
\def\faMemory{\textbf{[RAM]}}
\def\faServer{\textbf{[DC]}}
\def\faGlobe{\ensuremath{\diamond}}
\def\faStar{\ensuremath{\bigstar}}
\def\faMicrochip{\textbf{[$\mu$P]}}
\def\faLayerGroup{\textbf{[LAY]}}
\def\faCubes{\ensuremath{\boxtimes}}
\usepackage{amssymb}
\usepackage{pifont}
\usepackage{tcolorbox}
\tcbuselibrary{skins,breakable}
\usepackage{caption}
\captionsetup{font=small, skip=4pt}
\usepackage{float}
\usepackage{needspace}
\setlength{\textfloatsep}{8pt plus 2pt minus 2pt}
\setlength{\floatsep}{6pt plus 2pt minus 2pt}
\setlength{\intextsep}{6pt plus 2pt minus 2pt}
\setlength{\abovecaptionskip}{4pt}
\setlength{\belowcaptionskip}{2pt}

% ============ 自定义命令 ============
\newcommand{\rmark}{\textcolor{riskred}{\ding{108}}}
\newcommand{\amark}{\textcolor{riskamber}{\ding{108}}}
\newcommand{\gmark}{\textcolor{riskgreen}{\ding{108}}}
\newcommand{\starfive}{$\bigstar\bigstar\bigstar\bigstar\bigstar$}
\newcommand{\starfour}{$\bigstar\bigstar\bigstar\bigstar$}
\newcommand{\starthree}{$\bigstar\bigstar\bigstar$}
% 资料来源注：使用楷体（STKaiti）而非斜体，避免中文无斜体字形导致的字体回退警告
% \mdseries 避免继承周围 \bfseries 导致 STKaiti 粗体缺失警告
% \sloppy + \hbadness 允许长文件名/URL 在必要时断行
% 局部重定义 \textbf → \heiti：STKaiti 无粗体，用黑体替代强调
\newcommand{\sourcenote}[1]{\par{\scriptsize\mdseries\color{mutedgrey}\kaishu\sloppy\hbadness=9999\def\textbf##1{{\heiti ##1}} 资料来源：#1}\vspace{0.2em}}
\newcommand{\qualitynote}[1]{\par{\scriptsize\mdseries\color{mutedgrey}\sloppy\def\textbf##1{{\heiti ##1}} 数据质量：#1}\vspace{0.2em}}
\newcommand{\evidencetag}[1]{%
  \tikz[baseline=(x.base)]\node[
    xshift=0pt,
    inner xsep=3.5pt,
    inner ysep=1.2pt,
    rounded corners=1.5pt,
    fill=panelgrey,
    draw=rulegrey,
    text=steelgrey,
    font=\scriptsize
  ] (x) {#1};%
}
\newcommand{\stance}[2]{\textbf{\textcolor{#1}{#2}}}

% 信息框
\newtcolorbox{keyinsight}[1][]{
  colback=lightgrey, colframe=navy,
  fonttitle=\bfseries, title={#1}, breakable,
  arc=2pt, boxrule=0.5pt, left=8pt, right=8pt, top=6pt, bottom=6pt
}
\newtcolorbox{riskbox}[1][]{
  colback=red!5, colframe=riskred,
  fonttitle=\bfseries, title={#1}, breakable,
  arc=2pt, boxrule=0.5pt, left=8pt, right=8pt, top=6pt, bottom=6pt
}
\newtcolorbox{dashboardbox}[1][]{
  enhanced, breakable,
  colback=white, colframe=navy,
  boxrule=0.7pt, arc=2pt,
  fonttitle=\bfseries\color{white},
  coltitle=white, colbacktitle=navy,
  title={#1},
  left=9pt, right=9pt, top=7pt, bottom=7pt
}
\newtcolorbox{houseviewbox}[1][]{
  enhanced, breakable,
  colback=lightgrey, colframe=deepnavy,
  borderline west={3pt}{0pt}{accentblue},
  boxrule=0.4pt, arc=2pt,
  fonttitle=\bfseries, title={#1},
  left=9pt, right=9pt, top=7pt, bottom=7pt
}
\newtcolorbox{sourcequalitybox}[1][]{
  enhanced, breakable,
  colback=yellow!6, colframe=riskamber,
  boxrule=0.45pt, arc=2pt,
  fonttitle=\bfseries, title={#1},
  left=9pt, right=9pt, top=7pt, bottom=7pt
}
% exhibitbox：表/图外框。left=7pt+right=7pt=14pt 内边距，before upper 内的 \linewidth 已扣除该内边距。
% 约定：\exhibitboxwidth = 真实内边距宽（tcolorbox 已扣好 padding 的 \linewidth）。
% 由于 tabularx/longtable 宽度参数的语义是「外部总宽（含两侧 \tabcolsep）」，内部列实际可用 = W - 2\tabcolsep。
% 所以 tabularx 直接传 \exhibitboxwidth（总宽 == 内边距宽），而其他用途（resizebox、tikz、p{...}）用 \exhibitwidth = \exhibitboxwidth - 2\tabcolsep。
\newlength{\exhibitboxwidth}
\newlength{\exhibitwidth}
\newtcolorbox{exhibitbox}[1][]{
  enhanced, breakable,
  colback=white, colframe=rulegrey,
  borderline north={1.2pt}{0pt}{navy},
  boxrule=0.35pt, arc=1pt,
  fonttitle=\bfseries\color{navy}, title={#1},
  left=7pt, right=7pt, top=6pt, bottom=6pt,
  before upper={\global\exhibitboxwidth=\linewidth\relax
    \global\exhibitwidth=\dimexpr\linewidth-2\tabcolsep\relax},
}
\newtcolorbox{disclosurebox}[1][]{
  enhanced, breakable,
  colback=lightgrey, colframe=steelgrey,
  boxrule=0.35pt, arc=1pt,
  fonttitle=\bfseries, title={#1},
  left=8pt, right=8pt, top=6pt, bottom=6pt
}

% ============ First-page dashboard cover ============
\newcommand{\astockcover}{%
\begin{titlepage}
\thispagestyle{empty}
\begin{tikzpicture}[remember picture, overlay]
  \fill[deepnavy] (current page.north west) rectangle ([yshift=-2.55cm]current page.north east);
  \fill[accentblue] ([yshift=-2.55cm]current page.north west) rectangle ([yshift=-2.66cm]current page.north east);
  \fill[deepnavy] (current page.south west) rectangle ([yshift=1.2cm]current page.south east);
\end{tikzpicture}
\vspace*{0.62cm}
{\color{white}\sffamily\bfseries\small \reportkicker}\\[0.18cm]
{\color{steelgrey}\sffamily\scriptsize \reportscope \quad|\quad \reporttype \quad|\quad \reportdate}\\[1.0cm]
{\fontsize{25}{30}\selectfont\bfseries\color{navy}\reporttitle}\\[0.25cm]
{\fontsize{13}{17}\selectfont\color{steelgrey}\reportsubtitle}\\[0.75cm]
\begin{dashboardbox}[决策摘要]
\begin{tabularx}{\textwidth}{L{3.0cm}X L{3.0cm}X}
\textbf{报告类型} & \reporttype & \textbf{数据截至} & \reportdatacutoff \\
\textbf{覆盖范围} & \reportscope & \textbf{作者} & \reportauthor \\
\end{tabularx}
\vspace{0.45em}
\begin{houseviewbox}[本机构观点]
\reporthouseview
\end{houseviewbox}
\begin{sourcequalitybox}[证据质量]
\reportquality
\end{sourcequalitybox}
\end{dashboardbox}
\vfill
{\scriptsize\color{steelgrey}\reportdisclaimer}
\end{titlepage}
}
.
% Design benchmark: JPM first-page dashboard, BlackRock house-view framing,
% Vanguard scenario discipline, and AB/JPM chartbook exhibit structure.

\providecommand{\reporttitle}{AStock Research Report}
\providecommand{\reportsubtitle}{House View, Evidence, Valuation and Risk}
\providecommand{\reportkicker}{INSTITUTIONAL EQUITY RESEARCH}
\providecommand{\reportscope}{CHINA A-SHARES}
\providecommand{\reportdate}{\today}
\providecommand{\reportdatacutoff}{Data cutoff not specified}
\providecommand{\reporttype}{Research note}
\providecommand{\reportauthor}{AStock Research Agent}
\providecommand{\reporthouseview}{State the house view in one decision-useful paragraph.}
\providecommand{\reportquality}{Source quality not specified.}
\providecommand{\reportdisclaimer}{本报告基于公开资料整理，不构成任何证券买卖建议。}

% ============ 页面布局 ============
\usepackage[top=1.8cm, bottom=1.9cm, left=1.8cm, right=1.8cm]{geometry}
\setlength{\parskip}{0.3em}
\setlength{\parindent}{2em}
\renewcommand{\baselinestretch}{1.12}
\raggedbottom

% ============ Fonts ============
% Prefer the installed Noto CJK family when available (for Linux builds).
% Otherwise retain the original MacTeX typography. This avoids platform
% probing via shell escape and keeps the source portable across both engines.
\IfFontExistsTF{Noto Serif CJK SC}{
  \setCJKmainfont{Noto Serif CJK SC}
  \setCJKsansfont{Noto Sans CJK SC}
  \setCJKmonofont{Noto Sans CJK SC}
  \setmainfont{Liberation Serif}
  \newcommand{\kaishu}{\CJKfontspec{Noto Serif CJK SC}}
  \newcommand{\heiti}{\CJKfontspec{Noto Sans CJK SC}}
  \newcommand{\songti}{\CJKfontspec{Noto Serif CJK SC}}
}{
  \setCJKmainfont[BoldFont={Heiti SC}]{STSong}
  \setCJKsansfont{Heiti SC}
  \setCJKmonofont{Heiti SC}
  \setmainfont{Times New Roman}
  \newcommand{\kaishu}{\CJKfontspec{STKaiti}}
  \newcommand{\heiti}{\CJKfontspec{Heiti SC}}
  \newcommand{\songti}{\CJKfontspec{STSong}}
}

% ============ 颜色（IB Navy/Grey palette） ============
\usepackage{xcolor}
\definecolor{navy}{HTML}{003366}
\definecolor{deepnavy}{HTML}{0B1F33}
\definecolor{steelgrey}{HTML}{4A5568}
\definecolor{mutedgrey}{HTML}{718096}
\definecolor{lightgrey}{HTML}{F7FAFC}
\definecolor{panelgrey}{HTML}{EDF2F7}
\definecolor{riskred}{HTML}{C53030}
\definecolor{riskamber}{HTML}{D69E2E}
\definecolor{riskgreen}{HTML}{38A169}
\definecolor{riskblue}{HTML}{2B6CB0}
\definecolor{accentblue}{HTML}{2B6CB0}
\definecolor{nvgreen}{HTML}{76B900}
\definecolor{rulegrey}{HTML}{CBD5E0}

% ============ 表格 ============
\usepackage{booktabs}
\usepackage{longtable}
\usepackage{tabularx}
\usepackage{multirow}
\usepackage{array}
\usepackage{makecell}
\usepackage{colortbl}
\renewcommand{\arraystretch}{1.12}
\newcolumntype{Y}{>{\centering\arraybackslash}X}
\newcolumntype{L}[1]{>{\raggedright\arraybackslash}p{#1}}
\newcolumntype{C}[1]{>{\centering\arraybackslash}p{#1}}
\newcolumntype{R}[1]{>{\raggedleft\arraybackslash}p{#1}}

% ============ 图形 ============
\usepackage{graphicx}
\usepackage{tikz}
\usetikzlibrary{shapes,arrows.meta,positioning,calc,backgrounds,fit,decorations.pathreplacing}
\usepackage{pgfplots}
\usepgfplotslibrary{polar}
\pgfplotsset{compat=1.18}
\usepackage[edges]{forest}

% ============ 页眉页脚（报告标题和日期需在main.tex中覆盖） ============
\usepackage{fancyhdr}
\pagestyle{fancy}
\fancyhf{}
\fancyhead[L]{\scriptsize\color{steelgrey} \reporttitle}
\fancyhead[R]{\scriptsize\color{steelgrey} \reportdate}
\fancyfoot[C]{\small\color{steelgrey} \thepage}
\fancyfoot[R]{\tiny\color{steelgrey} \reportdisclaimer}
\renewcommand{\headrulewidth}{0.4pt}
\renewcommand{\footrulewidth}{0.2pt}

% ============ 章节格式 ============
\usepackage{titlesec}
\titleformat{\chapter}[block]
  {\normalfont\LARGE\bfseries\color{navy}}
  {\thechapter}{0.9em}{}
\titlespacing*{\chapter}{0pt}{0.4em}{1.0em}
\titleformat{\section}
  {\normalfont\large\bfseries\color{navy}}
  {\thesection}{1em}{}
\titleformat{\subsection}
  {\normalfont\normalsize\bfseries\color{steelgrey}}
  {\thesubsection}{1em}{}

% ============ 超链接 ============
\usepackage{hyperref}
\hypersetup{
  colorlinks=true,
  linkcolor=navy,
  urlcolor=navy,
  citecolor=navy,
}

% ============ 其他工具 ============
\usepackage{enumitem}
\setlist{nosep, leftmargin=2em}
\IfFileExists{fontawesome5.sty}{\usepackage{fontawesome5}}{}
% ============ Class B 全局兜底（xeCJK + 缺字 fallback 阻断）：
% 当 fontawesome5 因 Times New Roman 缺字 / \FA 命名空间失败时，
% 禁止输出 aExclamationTriangle / aCloud 裸文本，一律降级为 amssymb 原生符号。
% 注：amssymb 在 L112 已加载，此处不重复。所有 \def 在本段落。
\def\faExclamationTriangle{\ensuremath{\textcolor{riskred}{\blacktriangle}}\kern-0.25em}
\def\faCloud{\ensuremath{\clubsuit}\kern-0.1em}
\def\faYenSign{\ensuremath{\yen}}
\def\faIndustry{\ensuremath{\S}}
\def\faCogs{\ensuremath{\ast}}
\def\faFlagCheckered{\ensuremath{\blacktriangleright}}
\def\faMemory{\textbf{[RAM]}}
\def\faServer{\textbf{[DC]}}
\def\faGlobe{\ensuremath{\diamond}}
\def\faStar{\ensuremath{\bigstar}}
\def\faMicrochip{\textbf{[$\mu$P]}}
\def\faLayerGroup{\textbf{[LAY]}}
\def\faCubes{\ensuremath{\boxtimes}}
\usepackage{amssymb}
\usepackage{pifont}
\usepackage{tcolorbox}
\tcbuselibrary{skins,breakable}
\usepackage{caption}
\captionsetup{font=small, skip=4pt}
\usepackage{float}
\usepackage{needspace}
\setlength{\textfloatsep}{8pt plus 2pt minus 2pt}
\setlength{\floatsep}{6pt plus 2pt minus 2pt}
\setlength{\intextsep}{6pt plus 2pt minus 2pt}
\setlength{\abovecaptionskip}{4pt}
\setlength{\belowcaptionskip}{2pt}

% ============ 自定义命令 ============
\newcommand{\rmark}{\textcolor{riskred}{\ding{108}}}
\newcommand{\amark}{\textcolor{riskamber}{\ding{108}}}
\newcommand{\gmark}{\textcolor{riskgreen}{\ding{108}}}
\newcommand{\starfive}{$\bigstar\bigstar\bigstar\bigstar\bigstar$}
\newcommand{\starfour}{$\bigstar\bigstar\bigstar\bigstar$}
\newcommand{\starthree}{$\bigstar\bigstar\bigstar$}
% 资料来源注：使用楷体（STKaiti）而非斜体，避免中文无斜体字形导致的字体回退警告
% \mdseries 避免继承周围 \bfseries 导致 STKaiti 粗体缺失警告
% \sloppy + \hbadness 允许长文件名/URL 在必要时断行
% 局部重定义 \textbf → \heiti：STKaiti 无粗体，用黑体替代强调
\newcommand{\sourcenote}[1]{\par{\scriptsize\mdseries\color{mutedgrey}\kaishu\sloppy\hbadness=9999\def\textbf##1{{\heiti ##1}} 资料来源：#1}\vspace{0.2em}}
\newcommand{\qualitynote}[1]{\par{\scriptsize\mdseries\color{mutedgrey}\sloppy\def\textbf##1{{\heiti ##1}} 数据质量：#1}\vspace{0.2em}}
\newcommand{\evidencetag}[1]{%
  \tikz[baseline=(x.base)]\node[
    xshift=0pt,
    inner xsep=3.5pt,
    inner ysep=1.2pt,
    rounded corners=1.5pt,
    fill=panelgrey,
    draw=rulegrey,
    text=steelgrey,
    font=\scriptsize
  ] (x) {#1};%
}
\newcommand{\stance}[2]{\textbf{\textcolor{#1}{#2}}}

% 信息框
\newtcolorbox{keyinsight}[1][]{
  colback=lightgrey, colframe=navy,
  fonttitle=\bfseries, title={#1}, breakable,
  arc=2pt, boxrule=0.5pt, left=8pt, right=8pt, top=6pt, bottom=6pt
}
\newtcolorbox{riskbox}[1][]{
  colback=red!5, colframe=riskred,
  fonttitle=\bfseries, title={#1}, breakable,
  arc=2pt, boxrule=0.5pt, left=8pt, right=8pt, top=6pt, bottom=6pt
}
\newtcolorbox{dashboardbox}[1][]{
  enhanced, breakable,
  colback=white, colframe=navy,
  boxrule=0.7pt, arc=2pt,
  fonttitle=\bfseries\color{white},
  coltitle=white, colbacktitle=navy,
  title={#1},
  left=9pt, right=9pt, top=7pt, bottom=7pt
}
\newtcolorbox{houseviewbox}[1][]{
  enhanced, breakable,
  colback=lightgrey, colframe=deepnavy,
  borderline west={3pt}{0pt}{accentblue},
  boxrule=0.4pt, arc=2pt,
  fonttitle=\bfseries, title={#1},
  left=9pt, right=9pt, top=7pt, bottom=7pt
}
\newtcolorbox{sourcequalitybox}[1][]{
  enhanced, breakable,
  colback=yellow!6, colframe=riskamber,
  boxrule=0.45pt, arc=2pt,
  fonttitle=\bfseries, title={#1},
  left=9pt, right=9pt, top=7pt, bottom=7pt
}
% exhibitbox：表/图外框。left=7pt+right=7pt=14pt 内边距，before upper 内的 \linewidth 已扣除该内边距。
% 约定：\exhibitboxwidth = 真实内边距宽（tcolorbox 已扣好 padding 的 \linewidth）。
% 由于 tabularx/longtable 宽度参数的语义是「外部总宽（含两侧 \tabcolsep）」，内部列实际可用 = W - 2\tabcolsep。
% 所以 tabularx 直接传 \exhibitboxwidth（总宽 == 内边距宽），而其他用途（resizebox、tikz、p{...}）用 \exhibitwidth = \exhibitboxwidth - 2\tabcolsep。
\newlength{\exhibitboxwidth}
\newlength{\exhibitwidth}
\newtcolorbox{exhibitbox}[1][]{
  enhanced, breakable,
  colback=white, colframe=rulegrey,
  borderline north={1.2pt}{0pt}{navy},
  boxrule=0.35pt, arc=1pt,
  fonttitle=\bfseries\color{navy}, title={#1},
  left=7pt, right=7pt, top=6pt, bottom=6pt,
  before upper={\global\exhibitboxwidth=\linewidth\relax
    \global\exhibitwidth=\dimexpr\linewidth-2\tabcolsep\relax},
}
\newtcolorbox{disclosurebox}[1][]{
  enhanced, breakable,
  colback=lightgrey, colframe=steelgrey,
  boxrule=0.35pt, arc=1pt,
  fonttitle=\bfseries, title={#1},
  left=8pt, right=8pt, top=6pt, bottom=6pt
}

% ============ First-page dashboard cover ============
\newcommand{\astockcover}{%
\begin{titlepage}
\thispagestyle{empty}
\begin{tikzpicture}[remember picture, overlay]
  \fill[deepnavy] (current page.north west) rectangle ([yshift=-2.55cm]current page.north east);
  \fill[accentblue] ([yshift=-2.55cm]current page.north west) rectangle ([yshift=-2.66cm]current page.north east);
  \fill[deepnavy] (current page.south west) rectangle ([yshift=1.2cm]current page.south east);
\end{tikzpicture}
\vspace*{0.62cm}
{\color{white}\sffamily\bfseries\small \reportkicker}\\[0.18cm]
{\color{steelgrey}\sffamily\scriptsize \reportscope \quad|\quad \reporttype \quad|\quad \reportdate}\\[1.0cm]
{\fontsize{25}{30}\selectfont\bfseries\color{navy}\reporttitle}\\[0.25cm]
{\fontsize{13}{17}\selectfont\color{steelgrey}\reportsubtitle}\\[0.75cm]
\begin{dashboardbox}[决策摘要]
\begin{tabularx}{\textwidth}{L{3.0cm}X L{3.0cm}X}
\textbf{报告类型} & \reporttype & \textbf{数据截至} & \reportdatacutoff \\
\textbf{覆盖范围} & \reportscope & \textbf{作者} & \reportauthor \\
\end{tabularx}
\vspace{0.45em}
\begin{houseviewbox}[本机构观点]
\reporthouseview
\end{houseviewbox}
\begin{sourcequalitybox}[证据质量]
\reportquality
\end{sourcequalitybox}
\end{dashboardbox}
\vfill
{\scriptsize\color{steelgrey}\reportdisclaimer}
\end{titlepage}
}
.
% Design benchmark: JPM first-page dashboard, BlackRock house-view framing,
% Vanguard scenario discipline, and AB/JPM chartbook exhibit structure.

\providecommand{\reporttitle}{AStock Research Report}
\providecommand{\reportsubtitle}{House View, Evidence, Valuation and Risk}
\providecommand{\reportkicker}{INSTITUTIONAL EQUITY RESEARCH}
\providecommand{\reportscope}{CHINA A-SHARES}
\providecommand{\reportdate}{\today}
\providecommand{\reportdatacutoff}{Data cutoff not specified}
\providecommand{\reporttype}{Research note}
\providecommand{\reportauthor}{AStock Research Agent}
\providecommand{\reporthouseview}{State the house view in one decision-useful paragraph.}
\providecommand{\reportquality}{Source quality not specified.}
\providecommand{\reportdisclaimer}{本报告基于公开资料整理，不构成任何证券买卖建议。}

% ============ 页面布局 ============
\usepackage[top=1.8cm, bottom=1.9cm, left=1.8cm, right=1.8cm]{geometry}
\setlength{\parskip}{0.3em}
\setlength{\parindent}{2em}
\renewcommand{\baselinestretch}{1.12}
\raggedbottom

% ============ Fonts ============
% Prefer the installed Noto CJK family when available (for Linux builds).
% Otherwise retain the original MacTeX typography. This avoids platform
% probing via shell escape and keeps the source portable across both engines.
\IfFontExistsTF{Noto Serif CJK SC}{
  \setCJKmainfont{Noto Serif CJK SC}
  \setCJKsansfont{Noto Sans CJK SC}
  \setCJKmonofont{Noto Sans CJK SC}
  \setmainfont{Liberation Serif}
  \newcommand{\kaishu}{\CJKfontspec{Noto Serif CJK SC}}
  \newcommand{\heiti}{\CJKfontspec{Noto Sans CJK SC}}
  \newcommand{\songti}{\CJKfontspec{Noto Serif CJK SC}}
}{
  \setCJKmainfont[BoldFont={Heiti SC}]{STSong}
  \setCJKsansfont{Heiti SC}
  \setCJKmonofont{Heiti SC}
  \setmainfont{Times New Roman}
  \newcommand{\kaishu}{\CJKfontspec{STKaiti}}
  \newcommand{\heiti}{\CJKfontspec{Heiti SC}}
  \newcommand{\songti}{\CJKfontspec{STSong}}
}

% ============ 颜色（IB Navy/Grey palette） ============
\usepackage{xcolor}
\definecolor{navy}{HTML}{003366}
\definecolor{deepnavy}{HTML}{0B1F33}
\definecolor{steelgrey}{HTML}{4A5568}
\definecolor{mutedgrey}{HTML}{718096}
\definecolor{lightgrey}{HTML}{F7FAFC}
\definecolor{panelgrey}{HTML}{EDF2F7}
\definecolor{riskred}{HTML}{C53030}
\definecolor{riskamber}{HTML}{D69E2E}
\definecolor{riskgreen}{HTML}{38A169}
\definecolor{riskblue}{HTML}{2B6CB0}
\definecolor{accentblue}{HTML}{2B6CB0}
\definecolor{nvgreen}{HTML}{76B900}
\definecolor{rulegrey}{HTML}{CBD5E0}

% ============ 表格 ============
\usepackage{booktabs}
\usepackage{longtable}
\usepackage{tabularx}
\usepackage{multirow}
\usepackage{array}
\usepackage{makecell}
\usepackage{colortbl}
\renewcommand{\arraystretch}{1.12}
\newcolumntype{Y}{>{\centering\arraybackslash}X}
\newcolumntype{L}[1]{>{\raggedright\arraybackslash}p{#1}}
\newcolumntype{C}[1]{>{\centering\arraybackslash}p{#1}}
\newcolumntype{R}[1]{>{\raggedleft\arraybackslash}p{#1}}

% ============ 图形 ============
\usepackage{graphicx}
\usepackage{tikz}
\usetikzlibrary{shapes,arrows.meta,positioning,calc,backgrounds,fit,decorations.pathreplacing}
\usepackage{pgfplots}
\usepgfplotslibrary{polar}
\pgfplotsset{compat=1.18}
\usepackage[edges]{forest}

% ============ 页眉页脚（报告标题和日期需在main.tex中覆盖） ============
\usepackage{fancyhdr}
\pagestyle{fancy}
\fancyhf{}
\fancyhead[L]{\scriptsize\color{steelgrey} \reporttitle}
\fancyhead[R]{\scriptsize\color{steelgrey} \reportdate}
\fancyfoot[C]{\small\color{steelgrey} \thepage}
\fancyfoot[R]{\tiny\color{steelgrey} \reportdisclaimer}
\renewcommand{\headrulewidth}{0.4pt}
\renewcommand{\footrulewidth}{0.2pt}

% ============ 章节格式 ============
\usepackage{titlesec}
\titleformat{\chapter}[block]
  {\normalfont\LARGE\bfseries\color{navy}}
  {\thechapter}{0.9em}{}
\titlespacing*{\chapter}{0pt}{0.4em}{1.0em}
\titleformat{\section}
  {\normalfont\large\bfseries\color{navy}}
  {\thesection}{1em}{}
\titleformat{\subsection}
  {\normalfont\normalsize\bfseries\color{steelgrey}}
  {\thesubsection}{1em}{}

% ============ 超链接 ============
\usepackage{hyperref}
\hypersetup{
  colorlinks=true,
  linkcolor=navy,
  urlcolor=navy,
  citecolor=navy,
}

% ============ 其他工具 ============
\usepackage{enumitem}
\setlist{nosep, leftmargin=2em}
\IfFileExists{fontawesome5.sty}{\usepackage{fontawesome5}}{}
% ============ Class B 全局兜底（xeCJK + 缺字 fallback 阻断）：
% 当 fontawesome5 因 Times New Roman 缺字 / \FA 命名空间失败时，
% 禁止输出 aExclamationTriangle / aCloud 裸文本，一律降级为 amssymb 原生符号。
% 注：amssymb 在 L112 已加载，此处不重复。所有 \def 在本段落。
\def\faExclamationTriangle{\ensuremath{\textcolor{riskred}{\blacktriangle}}\kern-0.25em}
\def\faCloud{\ensuremath{\clubsuit}\kern-0.1em}
\def\faYenSign{\ensuremath{\yen}}
\def\faIndustry{\ensuremath{\S}}
\def\faCogs{\ensuremath{\ast}}
\def\faFlagCheckered{\ensuremath{\blacktriangleright}}
\def\faMemory{\textbf{[RAM]}}
\def\faServer{\textbf{[DC]}}
\def\faGlobe{\ensuremath{\diamond}}
\def\faStar{\ensuremath{\bigstar}}
\def\faMicrochip{\textbf{[$\mu$P]}}
\def\faLayerGroup{\textbf{[LAY]}}
\def\faCubes{\ensuremath{\boxtimes}}
\usepackage{amssymb}
\usepackage{pifont}
\usepackage{tcolorbox}
\tcbuselibrary{skins,breakable}
\usepackage{caption}
\captionsetup{font=small, skip=4pt}
\usepackage{float}
\usepackage{needspace}
\setlength{\textfloatsep}{8pt plus 2pt minus 2pt}
\setlength{\floatsep}{6pt plus 2pt minus 2pt}
\setlength{\intextsep}{6pt plus 2pt minus 2pt}
\setlength{\abovecaptionskip}{4pt}
\setlength{\belowcaptionskip}{2pt}

% ============ 自定义命令 ============
\newcommand{\rmark}{\textcolor{riskred}{\ding{108}}}
\newcommand{\amark}{\textcolor{riskamber}{\ding{108}}}
\newcommand{\gmark}{\textcolor{riskgreen}{\ding{108}}}
\newcommand{\starfive}{$\bigstar\bigstar\bigstar\bigstar\bigstar$}
\newcommand{\starfour}{$\bigstar\bigstar\bigstar\bigstar$}
\newcommand{\starthree}{$\bigstar\bigstar\bigstar$}
% 资料来源注：使用楷体（STKaiti）而非斜体，避免中文无斜体字形导致的字体回退警告
% \mdseries 避免继承周围 \bfseries 导致 STKaiti 粗体缺失警告
% \sloppy + \hbadness 允许长文件名/URL 在必要时断行
% 局部重定义 \textbf → \heiti：STKaiti 无粗体，用黑体替代强调
\newcommand{\sourcenote}[1]{\par{\scriptsize\mdseries\color{mutedgrey}\kaishu\sloppy\hbadness=9999\def\textbf##1{{\heiti ##1}} 资料来源：#1}\vspace{0.2em}}
\newcommand{\qualitynote}[1]{\par{\scriptsize\mdseries\color{mutedgrey}\sloppy\def\textbf##1{{\heiti ##1}} 数据质量：#1}\vspace{0.2em}}
\newcommand{\evidencetag}[1]{%
  \tikz[baseline=(x.base)]\node[
    xshift=0pt,
    inner xsep=3.5pt,
    inner ysep=1.2pt,
    rounded corners=1.5pt,
    fill=panelgrey,
    draw=rulegrey,
    text=steelgrey,
    font=\scriptsize
  ] (x) {#1};%
}
\newcommand{\stance}[2]{\textbf{\textcolor{#1}{#2}}}

% 信息框
\newtcolorbox{keyinsight}[1][]{
  colback=lightgrey, colframe=navy,
  fonttitle=\bfseries, title={#1}, breakable,
  arc=2pt, boxrule=0.5pt, left=8pt, right=8pt, top=6pt, bottom=6pt
}
\newtcolorbox{riskbox}[1][]{
  colback=red!5, colframe=riskred,
  fonttitle=\bfseries, title={#1}, breakable,
  arc=2pt, boxrule=0.5pt, left=8pt, right=8pt, top=6pt, bottom=6pt
}
\newtcolorbox{dashboardbox}[1][]{
  enhanced, breakable,
  colback=white, colframe=navy,
  boxrule=0.7pt, arc=2pt,
  fonttitle=\bfseries\color{white},
  coltitle=white, colbacktitle=navy,
  title={#1},
  left=9pt, right=9pt, top=7pt, bottom=7pt
}
\newtcolorbox{houseviewbox}[1][]{
  enhanced, breakable,
  colback=lightgrey, colframe=deepnavy,
  borderline west={3pt}{0pt}{accentblue},
  boxrule=0.4pt, arc=2pt,
  fonttitle=\bfseries, title={#1},
  left=9pt, right=9pt, top=7pt, bottom=7pt
}
\newtcolorbox{sourcequalitybox}[1][]{
  enhanced, breakable,
  colback=yellow!6, colframe=riskamber,
  boxrule=0.45pt, arc=2pt,
  fonttitle=\bfseries, title={#1},
  left=9pt, right=9pt, top=7pt, bottom=7pt
}
% exhibitbox：表/图外框。left=7pt+right=7pt=14pt 内边距，before upper 内的 \linewidth 已扣除该内边距。
% 约定：\exhibitboxwidth = 真实内边距宽（tcolorbox 已扣好 padding 的 \linewidth）。
% 由于 tabularx/longtable 宽度参数的语义是「外部总宽（含两侧 \tabcolsep）」，内部列实际可用 = W - 2\tabcolsep。
% 所以 tabularx 直接传 \exhibitboxwidth（总宽 == 内边距宽），而其他用途（resizebox、tikz、p{...}）用 \exhibitwidth = \exhibitboxwidth - 2\tabcolsep。
\newlength{\exhibitboxwidth}
\newlength{\exhibitwidth}
\newtcolorbox{exhibitbox}[1][]{
  enhanced, breakable,
  colback=white, colframe=rulegrey,
  borderline north={1.2pt}{0pt}{navy},
  boxrule=0.35pt, arc=1pt,
  fonttitle=\bfseries\color{navy}, title={#1},
  left=7pt, right=7pt, top=6pt, bottom=6pt,
  before upper={\global\exhibitboxwidth=\linewidth\relax
    \global\exhibitwidth=\dimexpr\linewidth-2\tabcolsep\relax},
}
\newtcolorbox{disclosurebox}[1][]{
  enhanced, breakable,
  colback=lightgrey, colframe=steelgrey,
  boxrule=0.35pt, arc=1pt,
  fonttitle=\bfseries, title={#1},
  left=8pt, right=8pt, top=6pt, bottom=6pt
}

% ============ First-page dashboard cover ============
\newcommand{\astockcover}{%
\begin{titlepage}
\thispagestyle{empty}
\begin{tikzpicture}[remember picture, overlay]
  \fill[deepnavy] (current page.north west) rectangle ([yshift=-2.55cm]current page.north east);
  \fill[accentblue] ([yshift=-2.55cm]current page.north west) rectangle ([yshift=-2.66cm]current page.north east);
  \fill[deepnavy] (current page.south west) rectangle ([yshift=1.2cm]current page.south east);
\end{tikzpicture}
\vspace*{0.62cm}
{\color{white}\sffamily\bfseries\small \reportkicker}\\[0.18cm]
{\color{steelgrey}\sffamily\scriptsize \reportscope \quad|\quad \reporttype \quad|\quad \reportdate}\\[1.0cm]
{\fontsize{25}{30}\selectfont\bfseries\color{navy}\reporttitle}\\[0.25cm]
{\fontsize{13}{17}\selectfont\color{steelgrey}\reportsubtitle}\\[0.75cm]
\begin{dashboardbox}[决策摘要]
\begin{tabularx}{\textwidth}{L{3.0cm}X L{3.0cm}X}
\textbf{报告类型} & \reporttype & \textbf{数据截至} & \reportdatacutoff \\
\textbf{覆盖范围} & \reportscope & \textbf{作者} & \reportauthor \\
\end{tabularx}
\vspace{0.45em}
\begin{houseviewbox}[本机构观点]
\reporthouseview
\end{houseviewbox}
\begin{sourcequalitybox}[证据质量]
\reportquality
\end{sourcequalitybox}
\end{dashboardbox}
\vfill
{\scriptsize\color{steelgrey}\reportdisclaimer}
\end{titlepage}
}

\setlength{\headheight}{20pt}
\AtBeginEnvironment{tabularx}{\setlength{\linewidth}{\dimexpr\linewidth-21.9pt\relax}}

\hypersetup{
  pdfauthor={\reportauthor},
  pdftitle={\reporttitle}
}

\begin{document}

\astockcover

\tableofcontents
\clearpage

\chapter{结论：复制投资委员会，不复制交易}
\begin{keyinsight}[一句话答案]
\textbf{有。}个人在中国市场能够使用的私募式方法，不是跟随某家私募的持仓，也不是小型杠杆收购（LBO）；而是把投资委员会的承保标准移植到公开市场：\textbf{先承保现金流与负债，再讨论增长；先写出可证伪的关键经营指标（KPI）与熊市情景，再决定是否进入研究候选池；只有证据增强才重新评估。} 在费用、账户权限、流动性、税费和个人约束均已核验的前提下，广泛分散的低成本核心工具与可为零的直接研究框架，是可比较的研究候选，而非对任何人的配置建议。
\end{keyinsight}

\section{本报告的决策结论}

全球另类资产管理人的披露呈现出多项机构结构与能力：控制权、公司经营干预、专属项目源、资本结构设计、基金或保险资金的期限匹配，以及研究和风控组织。把这些结构误读为一个选股秘诀，往往是个人投资者模仿失败的起点。官方披露显示，Blackstone、KKR、Apollo、Brookfield、Carlyle 与 TPG 的业务均横跨私募股权、信贷、实物资产、基金解决方案或保险相关资本；它们并不是只靠“低买高卖”的单一买断策略 [S01--S05, S07]。这些材料描述业务和流程，不构成对历史业绩成因或个人可复制回报的证明。

\begin{exhibitbox}[Exhibit 1：个人可用的结论地图]
\small
\begin{tabularx}{\exhibitboxwidth}{L{3.2cm}X L{3.0cm}}
\toprule
\textbf{问题} & \textbf{结论} & \textbf{实际含义} \\
\midrule
能否复制私募的收益？ & 不能把机构历史结果直接复制到个人账户。 & 内部收益率（IRR）、费率、杠杆、现金调用、估值频率、退出时点与流动性均不同。 \\
能否复制私募的过程？ & 可以借用其中的一部分承保和证伪纪律。 & 用公开披露、情景估值、关键经营指标（KPI）、债务压力测试和预先写好的失效条件。 \\
本框架的国内表达？ & “私募式公开市场承保 + 可为零的直接研究”。 & 直接股票、公募不动产投资信托基金（REIT）和事件研究都可以为零；不是必须持有某一类资产。 \\
最危险的误读？ & 把融资、私贷、高分派或并购故事当成私募优势。 & 默认零杠杆；高收益和低估值均不是安全证明。 \\
\bottomrule
\end{tabularx}
\sourcenote{S01--S07；本报告的分析归纳。}
\end{exhibitbox}

\section{三层迁移结论}

\begin{dashboardbox}[迁移优先级]
\begin{tabularx}{\textwidth}{L{2.4cm}X X}
\toprule
\textbf{层级} & \textbf{可用内容} & \textbf{对个人的约束} \\
\midrule
\textcolor{riskgreen}{高度可迁移} & 窄能力圈、现金流/债务承保、治理审查、熊/基/乐观情景、关键经营指标（KPI）、估值纪律、资金期限匹配、独立反方审查。 & 需要读披露和持续复核；不适合只靠行情或单一指标操作。 \\
\textcolor{riskamber}{部分可迁移} & 已由管理层公开启动的运营改善、正式公告驱动的事件研究、以项目现金流为中心的公募 REIT 分析。 & 个人只能观察和选择，不能推动经营；每个项目或事件须单独建档。 \\
\textcolor{riskred}{不可迁移 / 不应模仿} & 控制权、董事会席位、杠杆收购（LBO）、私募信贷发起、保险资产负债管理、私募二级份额、稳定借券空头和非公开信息。 & 这些不是“需要努力学习”的技能缺口，而是法律、结构、资源与风险承受能力的边界。 \\
\bottomrule
\end{tabularx}
\end{dashboardbox}

\section{本报告不会做什么}

本报告没有命名股票、交易型开放式指数基金（ETF）、不动产投资信托基金（REIT）或私募产品；没有使用实时价格、用户账户或风险偏好；也没有构造收益率目标。因此它只能作为研究流程的证据备忘录，而不是个性化配置或执行许可。机器可读的实施边界见\ \texttt{data/portfolio\_return\_model\_20260813.json}：所有市场敞口均为零，任何未来个券都必须先完成独立的公开信息、估值、流动性、账户与风险预算审查。

\begin{riskbox}[先避免错误问题]
“Blackstone / KKR 在买什么？”通常不是个人最有用的问题。更有用的问题是：\textbf{我是否能用公开资料写出一家企业或一个项目的现金流、债务、运营改善、估值与失效条件，而且在普通市场流动性下承担得起其熊市结果？} 如果答案是否定的，最专业的动作是暂不配置，而不是寻找更刺激的替代品。
\end{riskbox}


\chapter{全球头部私募到底在做什么}
\section{先纠正一个概念：参考管理人不是单一“收购策略”}

把本报告参考的七家管理人一概称为“私募”，容易掩盖它们之间的差异。Blackstone 的官方披露覆盖房地产、私募股权、信贷与保险、多资产、基础设施、生命科学和成长策略，并使用不同期限的资本载体 [S01]；KKR 覆盖私募股权、实物资产、信贷与流动策略、资本市场及保险相关业务 [S02]；Apollo 的核心语境则明显偏向发起、资产支持和多信用策略以及退休服务 [S03]。Brookfield 更强调实物资产的所有者兼经营者模式 [S04]；Carlyle 同时拥有全球私募、全球信贷与 AlpInvest 等基金解决方案 [S05]；TPG 则按资本、成长、影响力、信贷、房地产与市场解决方案等平台组织 [S07]。

\begin{exhibitbox}[Exhibit 2：七家机构披露的结构与本报告的研究转换]
\scriptsize
\begin{longtable}{L{2.0cm} L{3.4cm} L{4.0cm} L{4.0cm}}
\toprule
\textbf{机构} & \textbf{官方材料所描述的业务/方法} & \textbf{披露所反映的机构结构/能力} & \textbf{本报告的公开市场研究转换} \\
\midrule
\endfirsthead
\toprule
\textbf{机构} & \textbf{官方材料所描述的业务/方法} & \textbf{披露所反映的机构结构/能力} & \textbf{本报告的公开市场研究转换} \\
\midrule
\endhead
Blackstone & 多策略平台；私募、房地产、信贷/保险、基础设施和长期资本载体。 & 披露呈现多策略平台、长期资本，以及部分控制型私募策略中的控制/所有权。 & 购买前先写经营价值桥：收入、利润率、营运资本、维持性资本开支与债务服务；不把估值扩张当成唯一收益来源。 \\
KKR & 私募、实物资产、信贷/流动策略、资本市场与保险相关资本。 & 披露呈现投后运营资源、长期资金与资本结构；Global Atlantic 保险业务有资产负债现金流匹配语境。 & 先问债务期限、利息负担、抵押/担保和再融资，再讨论上行；投资期限必须匹配自己的现金需求。 \\
Apollo & 直投发起、资产支持、多信用、机会型信贷与退休服务。 & 披露呈现发起渠道、结构化资本解决方案、优先级与资产支持的承保语境，以及长期负债。 & 对股票也采用“债权人先问”的下行测试：企业的可持续现金是否能穿越压力期并覆盖义务？ \\
Brookfield & 基础设施、能源转型、房地产、私募与信贷；强调经营型实物资产。 & 管理人材料呈现运营业务、长期私募基金以及所有者兼经营者的经验。 & 对上市运营资产或公募不动产投资信托基金（REIT），研究合同/监管现金流、维护资本开支、杠杆、续约/出租率/业务量，而非只看分派率。 \\
Carlyle & 全球私募、全球信贷、基金一级/二级/共同投资与组合融资。 & 披露呈现全球网络和跨基金方案。 & 只跟踪与投资假设直接相连的少量披露关键经营指标（KPI）：订单、交付、利用率、入住率、预订或回款。 \\
Bain Capital & 私募及相邻资本策略；强调运营转型与价值创造。 & 管理人自述的运营转型方法与公司访问条件；并未在本来源中证明固定的动作清单或收益因果。 & 只有管理层已公开启动、可量化且能按季验证的改善，才可获得估值信用；不能为个人无权推动的改造提前付钱。 \\
TPG & 资本、成长、影响力、信贷、房地产、市场解决方案等平台。 & 披露呈现六个多策略平台；项目来源表述限于管理人自述。 & 用因果链做主题研究：主题 $\rightarrow$ 供需 $\rightarrow$ 收入/利润 $\rightarrow$ 估值 $\rightarrow$ 可证伪 KPI；主题本身不构成交易理由。 \\
\bottomrule
\end{longtable}
\sourcenote{S01--S07。前两列为官方材料的结构性描述或管理人自述；末列是本报告的分析转换。它们不构成各管理人收益的因果归因，也不证明个人可以复制回报。}
\end{exhibitbox}

\section{共同的机构能力组合：不是一个选股信号}

\begin{center}
\begin{tikzpicture}[node distance=0.55cm and 0.6cm, every node/.style={font=\small, align=center}]
\node[draw=navy, fill=lightgrey, rounded corners, minimum width=3.15cm, minimum height=1.15cm] (price) {价格与\\风险承保};
\node[draw=navy, fill=lightgrey, rounded corners, minimum width=3.15cm, minimum height=1.15cm, right=of price] (control) {控制权与\\投后运营};
\node[draw=navy, fill=lightgrey, rounded corners, minimum width=3.15cm, minimum height=1.15cm, right=of control] (source) {专属项目源与\\结构化条款};
\node[draw=navy, fill=lightgrey, rounded corners, minimum width=3.15cm, minimum height=1.15cm, below=of control] (capital) {匹配期限的资本\\与风险承受能力};
\node[draw=accentblue, fill=accentblue!12, rounded corners, minimum width=5.1cm, minimum height=1.1cm, below=of capital] (outcome) {机构流程的分析框架：\\资产选择、投后行动、资本结构与期限共同参与};
\draw[-{Stealth[length=2mm]}, thick, navy] (price.south) |- (outcome.west);
\draw[-{Stealth[length=2mm]}, thick, navy] (control.south) -- (capital.north);
\draw[-{Stealth[length=2mm]}, thick, navy] (source.south) |- (outcome.east);
\draw[-{Stealth[length=2mm]}, thick, navy] (capital.south) -- (outcome.north);
\end{tikzpicture}
\end{center}
\sourcenote{S01--S06；框图为跨机构流程归纳，不是历史业绩的因果拆解。}

其中，个人最容易观察到的是价格、财报和公开关键经营指标（KPI）；最难、也最不应假装拥有的是控制权、条款、发起渠道和低成本长期资本。这正是本报告后文把“承保”放在中心、把“模仿交易”放在禁区的原因。上述组合只是对披露信息的分析归纳，不是管理人回报的因果解释。

\section{为什么不能直接比较私募内部收益率（IRR）与个人账户收益}

私募基金的 IRR 受资本调用和分配时间、未上市资产估值、基金费用、项目杠杆、退出时点和流动性限制影响；个人账户更接近每天可见市值、连续交易成本和可变现金流的时间加权回报。两者都可能是有效的度量，但不是同一个问题。因此，管理规模、品牌、历史 IRR 或某一项目的传奇案例，不能被推导成中国个人账户的预期收益率。本报告只把管理人的官方材料用于回答“它们如何描述自己的业务和流程”，并把收益归因结论保留为未证实。

\begin{sourcequalitybox}[证据边界]
S01--S07 的级别足以支持机构结构与自述流程；其中 Bain 等运营价值创造表述尤其应被理解为管理人自述，而不是独立审计的因果证明。真正需要个人验证的不是“某家机构是否优秀”，而是未来某一个公开标的能否通过五票承保。
\end{sourcequalitybox}


\chapter{可迁移性：纪律、半迁移能力与禁区}
\section{高迁移：把“收购前委员会”搬到公开市场}

个人能复制的不是买断合同，而是更严苛的入场条件。公开市场可能提供较低的进入门槛和可选择不交易的自由，但费用、税费、流动性、生活现金需求、账户条件和行为偏差会显著改变实际体验；它们不构成天然优势。个人没有控制权，信息不对称仍在，且行情波动会持续考验纪律。因此，最有价值的私募式能力是下列六项：

\begin{enumerate}
  \item \textbf{能力圈收窄：}只研究少量、能持续跟踪的产业链；具体数量应由可用时间、披露频率和理解深度决定，而非追逐所有热门主题。
  \item \textbf{现金流先于叙事：}利润、经营现金流、应收、维持性资本开支、债务和再融资共同决定下行。
  \item \textbf{运营改善必须可见：}把提价、产品结构、降本、产能利用率、去杠杆或资产盘活写成可追踪的关键经营指标（KPI），而不是愿景。
  \item \textbf{价值桥而非单点估值：}把收益拆成盈利/现金流兑现、分红或分派、以及可能的估值变化；后者不能被默认。
  \item \textbf{分期拨款：}证据增强才增加承诺；股价下跌不是增加承诺的理由。
  \item \textbf{反方与失效：}在买入前写下最可能推翻自身结论的事实、日期和行动，而不是在下跌后临时解释。
\end{enumerate}

\section{部分迁移：能观察，不能推动}

运营改善、实物资产和事件驱动都可以成为公开市场研究对象，但其本质发生了变化。私募投资者可以任命管理层、调整预算或谈判融资；个人只能判断管理层已经做了什么、是否正在兑现，以及价值是否已被价格充分反映。对此，正确的估值处理是保守的：\textbf{没有披露、没有量化桥、没有关键经营指标（KPI）的“改善”，估值信用为零。}

\begin{exhibitbox}[Exhibit 3：迁移阶梯]
\small
\begin{tabularx}{\exhibitboxwidth}{L{2.7cm} L{3.0cm} X L{3.1cm}}
\toprule
\textbf{类别} & \textbf{机构原型} & \textbf{个人可做的动作} & \textbf{不可跨越的边界} \\
\midrule
高度可迁移 & KKR / Apollo 式承保 & 现金流、债务、再融资、估值、KPI、反方审查。 & 不能把“会做尽调”误认为拥有专属信息。 \\
高度可迁移 & 私募投资委员会（IC）/ 资本调用 & 以公开证据分阶段复核研究结论。 & 不因价格更低而自动加仓。 \\
部分可迁移 & Bain / Blackstone 式运营改善 & 选择已公开启动的经营改变并跟踪结果。 & 无法强迫管理层执行，不能把计划当利润。 \\
部分可迁移 & Brookfield 式实物资产 & 用项目现金流、维护资本开支和融资压力研究公募 REIT。 & 没有项目控制权，也不应把分派当无风险利息。 \\
部分可迁移 & 正式公告事件研究 & 围绕公告、业绩、资产处置或去杠杆建立可为零的研究档案。 & 传闻、未公开信息和“消息灵通”不属于研究优势。 \\
不迁移 & 杠杆收购（LBO）、控制权、私贷、保险资产负债管理（ALM） & 无。 & 法律结构、资金成本、管理资源与流动性容忍度均不同。 \\
\bottomrule
\end{tabularx}
\sourcenote{S01--S07；制度转换为本报告分析。}
\end{exhibitbox}

\section{明确的红线}

\begin{riskbox}[这不是个人的“待学习清单”]
以下内容不应被包装为个人可照搬的方法：控制权收购与私有化、董事会影响、杠杆收购（LBO）融资、困境债权人控制、私募直贷发起、基金二级份额交易、需要稳定借券的多空策略、并购套利基础设施，以及任何非公开信息。它们依赖法律地位、资金池、运营团队、合规系统、交易对手与承受长期锁定的能力。缺少这些条件时，模仿只会放大尾部风险。
\end{riskbox}

\section{个人真正可能拥有的比较优势}

个人没有机构的访问权，却可能在个人现金需求、账户条件、交易成本和行为约束都允许时，选择不强制配置、以可为零的研究范围等待公开证据，并在真正理解的领域长期积累。这些可能性不是天然优势，无法保证超额收益，却有助于避免“看不懂也参与”。若投资者没有兴趣或时间维护这套流程，合理的结果可以是直接研究卫星为零，并在费用、账户、流动性和税费均核验后考虑更简单的分散化方案。


\chapter{个人版核心：私募式公开市场承保}
\section{核心方法：五票全过，才有资格进入候选池}

个人版的私募式公开市场承保应像一个小型投资委员会，而不是一个因价格波动而不断改口的交易计划。每个候选对象必须同时通过五票；任何一票否决，最合理的结果都是“不买”或“继续研究”。这不是为了制造复杂度，而是为了让乐观故事必须同时经受财务、治理、估值和可执行性的约束。

\begin{exhibitbox}[Exhibit 4：五票承保卡]
\small
\begin{tabularx}{\exhibitboxwidth}{L{2.1cm} X L{3.55cm}}
\toprule
\textbf{投票} & \textbf{必须回答的问题} & \textbf{一票否决或降级信号} \\
\midrule
业务 & 公司如何赚钱？客户为何持续付费？核心假设如何被公开证伪？ & 只能用政策口号、价格走势或宏大主题解释盈利。 \\
现金与负债 & 利润是否变成现金？应收、库存、资本开支、净负债、担保和再融资是否经得起熊市？ & 现金长期背离利润且无可信解释；债务/担保足以击穿熊市假设。 \\
治理 & 审计、控股股东质押/减持、关联交易、激励和监管问询是否可解释？ & 非标审计、资金占用、重大治理恶化或无法定价的关联风险。 \\
估值 & 熊/基/乐观情景中，收入、利润率、资本开支和估值锚是什么？ & 仅靠历史市盈率（PE）、券商目标价或远期可触达市场规模（TAM）；基准情景需依赖多重扩张。 \\
验证与流动性 & 以与公司披露节奏匹配的观察期、哪个公开关键经营指标（KPI）证明或否定？正常市场下如何减仓？ & 没有公开证明点；假设可以随时按理想价格退出。 \\
\bottomrule
\end{tabularx}
\sourcenote{C03--C05；完整字段见\ \texttt{analysis/retail\_underwriting\_scorecard.md}。}
\end{exhibitbox}

\section{价值创造桥：只给已可见的动作估值信用}

定价、采购、渠道、产能、成本、营运资本、资产处置或管理激励，是企业可能采用的经营动作，并非本报告对某一管理人的固定价值创造清单。个人没有权利启动这些行动，所以应将“管理层可能改善”与“已披露、已量化、可追踪的改善”严格分开。一个合格的公开市场价值桥至少写清：

\begin{itemize}
  \item 机制：收入、销量、价格、产品组合、毛利率、费用率、营运资本或资本开支中的哪一项会变化；
  \item 基线：上一年或最近披露期的数值、口径和财务报表位置；
  \item 行动：谁已宣布什么行动，何时开始，是否需要监管、客户或融资条件；
  \item 关键经营指标（KPI）：下一份或未来若干份公开报告中可验证的指标；
  \item 估值信用：若 KPI 没有兑现，模型必须自动退回到无改善的保守情景。
\end{itemize}

这套规则尤其适合把“转型故事”从估值模型中剥离出来。研究者可以相信管理层，但模型不应提前把信念变成利润。

\section{从信用视角看股票：先问最坏情况下谁付账}

Apollo 的披露描述了信贷、资产支持与长期负债语境；KKR 对 Global Atlantic 保险业务披露了资产负债现金流匹配 [S02--S03]。这些机构语境提醒投资者，资本结构和期限不是配角。即使研究的是权益，也应先像债权人一样问：经营现金流是否覆盖利息和必要资本开支？债务到期是否集中？再融资是否依赖脆弱的市场环境？有无担保、质押、表外承诺或应收回款风险？企业的资产是否在下行周期仍有可变现或可经营价值？

这不是要求每家企业都低杠杆，而是要求杠杆风险可以被理解、建模并纳入熊市场景。高增长、高分派或低市净率不能替代这一步。

\section{分期资金投放：把“资本调用”改为公开证据调用}

\begin{center}
\begin{tikzpicture}[node distance=0.42cm, every node/.style={font=\small, align=center}]
\node[draw=navy, fill=lightgrey, rounded corners, minimum width=3.0cm, minimum height=1cm] (a) {初始承保\\五票通过};
\node[draw=navy, fill=lightgrey, rounded corners, minimum width=3.0cm, minimum height=1cm, right=of a] (b) {下一份报告\\验证核心指标};
\node[draw=navy, fill=lightgrey, rounded corners, minimum width=3.0cm, minimum height=1cm, right=of b] (c) {现金流/利润率\\验证机制};
\node[draw=accentblue, fill=accentblue!12, rounded corners, minimum width=3.0cm, minimum height=1cm, below=of b] (d) {重新承保：\\估值仍有余量？};
\draw[-{Stealth[length=2mm]}, thick, navy] (a) -- (b);
\draw[-{Stealth[length=2mm]}, thick, navy] (b) -- (c);
\draw[-{Stealth[length=2mm]}, thick, navy] (c.south) |- (d.east);
\draw[-{Stealth[length=2mm]}, thick, navy] (d.west) -| (a.south);
\end{tikzpicture}
\end{center}

每一步都应由新的公开证据触发；若证据没有增强，正确的结果是维持、缩小或归零，而不是因为价格下跌而自动摊薄成本。具体的初始仓位、单名上限和风险预算必须由个人资金需求和熊市损失承受能力决定，本报告不把示例比例伪装成适合所有人的配置。

\section{核心分散化与可为零的研究卫星}

私募基金带来的一个事实是，单一项目风险通常由组合和持续尽调稀释。个人无法得到同样的项目池，因此可把低成本、流动性充足的广泛公开市场工具作为基础分散化候选，把真正需要大量研究的直接股票、公募 REIT 或正式公告事件限制在可为零的研究卫星中。这样做的目的不是提高交易频率，而是避免把全部财务结果押在少数自认为理解的企业上。任何实际比例均须等账户、生活现金、风险承受能力、交易型开放式指数基金（ETF）穿透重叠和当前流动性完成后才能批准。


\chapter{中国市场如何改变仓位与退出设计}
\section{制度不是脚注：它决定了仓位设计}

把全球私募方法搬到中国个人账户时，最容易遗漏的是市场微观结构。机构可以依靠协议条款、长期资本、交叉资产流动性或专门交易团队管理退出；普通公开市场投资者需要面对交易制度、价格限制、停复牌、订单簿深度、账户权限和税费。于是，稳健的研究框架首先是一种\textbf{不会假设不存在的流动性}的设计。

\begin{exhibitbox}[Exhibit 5：中国个人版方法的制度适配]
\small
\begin{tabularx}{\exhibitboxwidth}{L{3.0cm} X L{4.1cm}}
\toprule
\textbf{约束} & \textbf{框架影响} & \textbf{研究动作} \\
\midrule
普通现货的结算 / 回转交易边界 & 普通股票应按交收前不得卖出的规则设计；有限例外必须逐只核验，不能把卖出设计成一定可当日执行的止损。 & 按非日内可处置的压力情景建仓；核验具体品种、现行登记结算和交易规则。 \\
涨跌幅、停牌和上市初期不连续性 & 价格限制不是全市场恒定数字：板块、基金类别、特殊情形和新上市阶段可不同；无涨跌幅也不等于一定有流动性。 & 排除或特别降级新上市、特殊处理、长期停牌、连续涨跌停或治理事件高发对象；限价与分期不等于保证成交。 \\
融资融券的准入和维持要求 & 融资并不等于可复制杠杆收购（LBO）的稳定资本结构，融券也不等于稳定可得的对冲。 & 默认零融资加仓、零裸空头；不让借款成为策略成立条件。 \\
交易型开放式指数基金（ETF）机制差异 & ETF 的二级市场、申赎和回转安排取决于类别、产品文件和当期规则；仅有某类 ETF 的机制不能外推到权益、跨境、商品或 REIT。 & 不把“可申赎”错当成二级市场必然的即时退出；查看基金合同、交易规则和流动性。 \\
公募不动产投资信托基金（REIT）产品结构 & 它是运营资产权益，不是现金替代物或固定利息工具。 & 对每个项目单独检查资产权利、经营指标、维护资本开支、融资、治理、估值、扩募/解禁和成交能力。 \\
\bottomrule
\end{tabularx}
\sourcenote{S08--S14。规则为框架级摘要；具体交易前须核验当日有效规则、产品文件与账户权限。}
\end{exhibitbox}

\section{为什么结算与回转交易边界改变“止损”思维}

上交所现行交易规则列出当日回转交易的特定例外；普通买入股票的一般处置应按交收前不得卖出的规则设计，并须尊重具体品种、现行登记结算和交易规则 [S09]。深交所的现行规则亦应在每次实施时对照确认 [S12]。因此，普通 A 股策略不应把“跌到某个价格就立即卖掉”当作完整的风险控制。价格阈值的正确作用是触发一次\textbf{重新承保}：事实是否改变？熊市假设是否被击穿？若需要减仓，当前价格限制、可用流动性和账户规则下的可行路径是什么？

这并非反对纪律，而是把纪律从机械订单改成更难但更真实的三件事：事前仓位足够小、事前流动性足够好、事前失效条件足够明确。若一项投资只有在完美盘口或同日卖出条件下才安全，它在进入时就不安全。

\section{交易型开放式指数基金（ETF）、融资融券与“伪机构化”}

不同 ETF 的二级市场、申购赎回与回转安排取决于产品类别、产品文件和当期规则；即使某类产品存在申赎机制，也不等同于个人在二级市场拥有无摩擦的即时退出，且不能将某类 ETF 的机制泛化至权益、跨境、商品或 REIT [S13]。类似地，融资融券受标的范围、期限、保证金、担保物、维持担保和券商合同等条件限制 [S14]。本框架把它们都视为需要单独核验的产品/账户属性，而非一个可拿来放大私募收益的按钮。

个人投资者最常见的伪机构化行为是：用融资增加看似“高确定性”的仓位、用高分派产品替代现金、或以为更多代码就意味着更分散。这三者都可能让表面上的复杂度掩盖同一个真实风险因子：资金紧张、利率敏感、地产/租金、特定产业需求、或流动性抽离。

\section{公募不动产投资信托基金（REIT）：可研究的实物资产代理，但不是私募基础设施}

Brookfield 的所有者兼经营者模式启发的是研究顺序，而不是可照搬的持有结构 [S04]。上交所的公开募集不动产 REIT 业务办法及其发售上市指南提供相关制度语境 [S10--S11]；该办法衔接基础设施 REIT 指引与商业不动产 REIT 试点语境，实际适用范围仍须由每个产品的底层资产类别、现行规则和产品文件确认。每一个项目和产品仍需单独回答：

\begin{itemize}
  \item 现金流来自出租、车流、能源业务量、价格机制还是其他权利？其可持续性与合同/监管基础是什么？
  \item 经营 KPI 是出租率、加权平均租约期限、回款、车流、利用率、价格政策还是客户集中度？
  \item 维护资本开支、债务到期、再融资、扩募、解禁、关联交易和运营商激励如何影响可分配现金？
  \item 价格相对项目价值和未来可分配现金的含义是什么？市场成交能否支持预期的减仓路径？
\end{itemize}

只有这些问题有现时、项目级的披露答案时，REIT 才是“部分可迁移”的实物资产研究对象。历史高分派率本身不是资产质量、流动性或安全边际的证明。

\begin{riskbox}[实施前的强制暂停]
没有账户权限、当前产品合同、税费、买卖价差/深度、正常日成交能力、投资者风险承受能力和经批准的组合风险预算时，任何品种都不得从本报告的 \texttt{q=0} 状态升级。本备忘录提供研究逻辑，不授予交易资格。
\end{riskbox}


\chapter{投后核验、再承保与失效机制}
\section{投后不是“长期持有”，而是持续再承保}

私募机构拥有投后团队、董事会信息和经营会；个人不具备这些渠道，所以投后流程必须更少、更公开、更可重复。最有效的做法不是全天盯盘，而是把每一项持有理由映射到一个公告、财报或项目运营指标，并在固定时点重新比较事实与原假设。

\begin{exhibitbox}[Exhibit 6：公开信息下的监测节奏]
\small
\begin{tabularx}{\exhibitboxwidth}{L{2.3cm} L{3.5cm} X L{3.4cm}}
\toprule
\textbf{频率} & \textbf{看什么} & \textbf{问题} & \textbf{可能行动} \\
\midrule
月度 & 重大公告、仓位、成交和经济因子重叠。 & 是否出现控制权、审计、融资、监管、停牌或流动性变化？ & 暂停新增；把异常提交红队复核。 \\
季报 / 半年报 & 核心关键经营指标（KPI）、现金流、应收/库存、利润率、债务与再融资。 & 价值桥是否按原口径兑现？熊市假设需要下调吗？ & 维持、减少或退出；不因价格更低而自动加仓。 \\
半年 & 估值、可替代机会和反方论证。 & 基准情景仍有足够余量吗？是否把已经发生的改善重复计价？ & 重新承保；去除无证据的上行。 \\
年度 & 能力圈与方法本身。 & 研究流程是否真的比简单分散化改善决策质量？ & 缩小或归零研究卫星，若过程不能持续执行。 \\
\bottomrule
\end{tabularx}
{\setlength{\emergencystretch}{2em}\sourcenote{本报告的流程框架；个券关键经营指标（KPI）必须来自其自身公开披露。}}
\end{exhibitbox}

\section{硬失效与软失效}

\begingroup
\setlength{\tabcolsep}{4pt}
\begin{tabularx}{\linewidth}{L{2.3cm}L{5.8cm}X}
\toprule
\textbf{失效类型} & \textbf{典型事实} & \textbf{正确反应} \\
\midrule
硬失效 & 财务造假、无法解释的非标审计、控制权或治理崩塌、核心业务/牌照永久受损、债务或再融资击穿熊市假设。 & 立即停止新增；核验事实并在实际流动性、交易规则和价格限制下规划退出。不要假设跌停日一定可以成交。 \\
软失效 & 连续报告期低于原熊市 KPI、预设催化期满仍无兑现、估值余量消失、与既有持仓出现未预料的高重叠。 & 重做价值桥与反方；若无法恢复可验证的上行，减至目标或退出。 \\
价格波动本身 & 一日或一段时间的上涨/下跌，但没有新增基本面事实。 & 触发下一交易日或下一个工作日的研究复核，而非自动把价格变化解释成证据。 \\
\bottomrule
\end{tabularx}
\endgroup

\section{公共关键经营指标（KPI）的选择}

KPI 必须是能把经营叙事和财务结果连接起来的桥，而不是越多越好。制造业可用交付、价格、产能利用率或产品组合；消费/服务可用同店、订单、回款或客户留存；资产运营可用入住率、租约、租户集中度、车流或利用率。每个对象应只保留少量最关键、且来自正式披露的指标。若公司没有披露足以验证核心假设的数据，结论不是去猜数据，而是降低研究等级或放弃。

\section{退出纪律与可选性}

在私募中，退出是预先规划、可能持续多年的过程。个人账户的对应物是：在建仓之前就知道什么事实会让自己不再愿意持有，并保留足够的流动性和未投入资本来处理不确定性。现金在本框架中不是“闲置”，而是避免被迫卖出、避免用杠杆修补错误、等待更好公开证据的可选性。它不保证收益，却是个人最能与机构长期资本相对应的风险管理工具。


\chapter{红队审查：这套框架在哪些情况下不值得用}
\section{核心反方：流程精致，不等于可以战胜市场}

这份框架最需要防范的，不是粗心，而是“把严谨感当成优势本身”。现金流模型、关键经营指标（KPI）和红队笔记可以减少明显错误，却不能消除估值风险、信息时滞、黑天鹅、治理突变、行业竞争或市场风险偏好变化。一个优秀企业仍可能在已被充分定价后带来差回报；一份公开财报也可能来得太晚，无法阻止损失。

\begin{riskbox}[红队结论]
即使公开资料充分、KPI 与估值模型完备，个人仍可能因信息时滞、治理突变、停牌、涨跌停或共同因子拥挤而无法按计划退出。如果投资者无法长期读披露、维护 KPI、承认失效、并忍受研究结论为“暂不买入”，那么这套方法未必比低成本广泛分散化更好。此时，直接研究卫星可以是零。框架的价值首先是改善风险控制和解释质量，而不是承诺超额收益（alpha），更不构成可复制私募超额收益的证据。
\end{riskbox}

\section{四个常见失败模式}

\begin{exhibitbox}[Exhibit 7：私募模仿失败的情景]
\small
\begin{tabularx}{\exhibitboxwidth}{L{3.2cm}X X}
\toprule
\textbf{失败模式} & \textbf{为何会发生} & \textbf{防护栏} \\
\midrule
把品牌当信号 & 看到顶级机构涉足某赛道，就跳过企业层面承保。 & 品牌只能提供研究问题，不能替代五票、估值和流动性审查。 \\
把高分派当安全 & 忽略维护资本开支、再融资、底层运营、估值和成交能力。 & 对 REIT/收益资产建立项目级可分配现金流桥，拒绝历史分派年化外推。 \\
把杠杆当确定性 & 低估保证金、价格限制、券源和再融资对个人账户的放大作用。 & 零融资加仓、零裸空头；不让任何借款成为回报实现前提。 \\
把复杂度当分散化 & 多个代码、ETF 或收益产品实际暴露于同一经济因子。 & 穿透到需求、利率、租金、商品、融资和客户集中度；真实重叠优先于代码数量。 \\
\bottomrule
\end{tabularx}
\sourcenote{S08--S14；本报告的反方推演。}
\end{exhibitbox}

\section{什么时候应当不用它}

以下情况应优先选择更简单的方案，而不是强行做“个人私募”：短期内有明确大额资金需求；没有稳定应急储备；无法承受正常市场的显著波动；没有时间或兴趣维护季度研究；不具备相应板块或产品的账户权限；或者发现自己的结论无法通过第二人的反方审查。方法的成熟表现不是永远参与，而是知道什么时候不参与。

\section{最终框架结论}

本报告的个人可用方法可以浓缩为一句话：\textbf{以私募标准做公开市场承保，用公开证据而不是控制权推动决策，把核心分散化与现金当作可选性，并把直接研究、REIT 和事件表达限制在经过独立验证、可为零的研究卫星。}

它既不需要神秘信息，也不依赖日内交易或杠杆。它要求的反而是较少的标的、更慢的行动、更严格的证据和更诚实的“不知道”。


\appendix
\chapter{一手来源、证据等级与使用边界}
\section{来源等级与可复核性}

本报告将源文件按全球机构和中国制度两类目录归档；URL、检索日期、文件 SHA-256、用途和局限记录在源注册表。材料以一手官方披露为主；并不把管理规模、媒体转述或搜索摘要当作估值或收益证据。

\begin{longtable}{L{1.05cm} L{3.6cm} L{3.0cm} L{5.15cm}}
\toprule
\textbf{ID} & \textbf{文件} & \textbf{用于支持} & \textbf{主要局限} \\
\midrule
\endfirsthead
\toprule
\textbf{ID} & \textbf{文件} & \textbf{用于支持} & \textbf{主要局限} \\
\midrule
\endhead
S01 & Blackstone 2025 Form 10-K & 多策略、长期/永续资本及基础设施等业务结构。 & 管理人文件描述其业务，不证明个人可复制其回报。 \\
S02 & KKR 2025 Form 10-K & 私募、实物资产、信贷/流动策略与保险相关资本结构。 & 无法把投后资源或资金期限移植至个人账户。 \\
S03 & Apollo 2025 Form 10-K & 直投发起、资产支持/多信用与退休服务。 & 支持信用承保类比，不支持零售进入私募直贷。 \\
S04 & Brookfield 2025 Annual Information Form & 经营型实物资产与基础设施相关语境。 & 个人没有所有者兼经营者的控制与融资能力。 \\
S05 & Carlyle 2025 Form 10-K & 全球私募、信贷与基金解决方案。 & 私募二级和共同投资不是公开市场策略。 \\
S06 & Bain Capital official 2025 report & 运营改善与价值创造的自述。 & 属管理人自述，非独立业绩归因。 \\
S07 & TPG 2025 Form 10-K & 六个平台的策略组织。 & 主题/平台不是选股或交易信号。 \\
S08 & 2014 年证监会融资融券教育页 & 历史监管/教育背景。 & 非现行操作规则，不作为当前制度结论依据。 \\
S09 / S12 & 沪深现行交易规则 & 回转交易、结算和市场结构的框架。 & 具体品种例外和当日适用规则须再次核验。 \\
S10 / S11 & 上交所公开募集不动产 REIT 业务办法及发售上市指南 & 公募不动产 REIT 的发行、上市与研究边界。 & 具体底层资产规则、产品文件、运营与估值仍须独立核验。 \\
S13 & 上交所 ETF 业务规则 & 申赎机制与二级市场处置的区分。 & 不替代某一 ETF 的产品文件和流动性检查。 \\
S14 & 上交所融资融券规则 & 交易所层面的实施边界。 & 不构成账户适当性判断。 \\
\bottomrule
\end{longtable}

\section{证据与结论的对应关系}

\begingroup
\setlength{\tabcolsep}{4pt}
\begin{tabularx}{\linewidth}{L{2.9cm}X X}
\toprule
\textbf{结论类型} & \textbf{可用证据} & \textbf{不可推出的内容} \\
\midrule
机构业务结构 & 官方 10-K、AIF、官方报告。 & 机构历史收益必然可复制。 \\
中国制度约束 & CSRC / 交易所现行规则。 & 某账户、某产品或某时点一定可执行。 \\
个人版承保流程 & 对机构方法与市场约束的分析转换。 & 某只证券的买卖结论。 \\
公募 REIT 研究方法 & 对应类别的规则、指南与未来项目披露。 & 历史分派就是未来安全收益。 \\
\bottomrule
\end{tabularx}
\endgroup

\noindent\begin{minipage}{0.98\linewidth}\sloppy
完整 claim-to-source 映射见\ \texttt{data/claim\_audit.json}；\allowbreak 未解决的证据缺口见\ \texttt{source\_exhaustion\_log.md}。
\end{minipage}


\clearpage
\thispagestyle{empty}
\vspace*{4cm}
\begin{disclosurebox}[重要披露]
\small
\reportdisclaimer\par
本备忘录通过“战略研究框架”审查，并不等同于实施批准。没有用户的资产负债、现金需求、风险承受能力、账户权限、税务条件、既有持仓、个券尽调、当前流动性或经批准的风险预算，因此任何市场交易敞口均为 \texttt{q=0}。本报告不承诺收益，也不把历史分派、私募内部收益率（IRR）、管理规模或品牌转化为未来回报预测。
\end{disclosurebox}

\end{document}
